\section{Dataset Process}
\label{app:dataset-process}

% 介绍原始数据的清洗工作
\subsection{Source Data and Data Cleaning}

We obtained the original data from \citet{sahinuc-etal-2026-reward} and retained 
three related-work and four review-utility dimensions, excluding 
\textit{verifiability extraction}. We removed duplicate, overlapping, malformed, 
and unreliable instances. RevUtil is substantially larger than
the related-work data; we therefore reproducibly
sampled 4,800 candidates per RevUtil dimension while approximately preserving
the original label distributions.

% 介绍三个工作：1.prompt 修改 2.蒸馏得到 cot 数据 3.构建 cot 和 label-only 2种数据
\subsection{Training and Test Data Construction}

\noindent\textbf{Zero-shot prompt construction.}
The original prompts included fixed few-shot examples containing both
rationales and labels, which could influence the model's reasoning and
confound comparisons across supervision settings. As illustrated in Figure
\hyperref[fig:original-prompt]{~\ref*{fig:original-prompt}} and Figure
\hyperref[fig:normalized-prompt]{~\ref*{fig:normalized-prompt}}, we
removed the example block and its corresponding system instruction while
leaving the query, criteria, scoring rubric, and evaluation instance
unchanged.

\noindent\textbf{Teacher rationale distillation.}
We use DeepSeek-V4-Pro to generate a rationale and score for each training
candidate under the configuration summarized in
Table~\ref{tab:distillation-configuration}. A trajectory is retained only if
it strictly follows the required output format, contains a score from the
task-specific label space, and assigns the same score as the gold label.

\begin{table}[H]
\centering
\small
\begin{tabular}{lp{0.58\columnwidth}}
\hline
\textbf{Configuration} & \textbf{Value} \\
\hline
Teacher model & DeepSeek-V4-Pro \\
Temperature & 0 \\
Top-$p$ & 1 \\
Maximum output tokens & 2,048 \\
Number of completions & 1 \\
Output format & \texttt{<reasoning>...</reasoning>} followed by
\texttt{<score>N</score>} \\
Acceptance criterion & Format valid and score matches gold label \\
\hline
\end{tabular}
\caption{Configuration used for teacher-rationale distillation.}
\label{tab:distillation-configuration}
\end{table}

\noindent\textbf{CoT and label-only data construction.}
The final training data contain two supervision formats. CoT data pair each
evaluation prompt with a teacher-generated rationale followed by the gold
score, whereas label-only data pair the prompt with the gold score alone.
Figures~\ref{fig:cot-example} and~\ref{fig:label-only-example} present the two
formats using the same RevUtil actionability instance.

% 展示数据统计
\subsection{Final Dataset Statistics}

Table~\ref{tab:dataset-process-statistics} reports the task format and final
numbers of training and test instances for all seven datasets, together with
their aggregate sizes.

\begin{table}[H]
\centering
\small
\begin{tabular}{p{0.43\columnwidth}crr}
\hline
\textbf{Task} & \textbf{Format} & \textbf{Train} & \textbf{Test} \\
\hline
Related Work--Coherence & Binary & 1,526 & 1,046 \\
Related Work--Positioning Check & Binary & 2,666 & 603 \\
Related Work--Positioning Type & Binary & 944 & 204 \\
RevUtil--Actionability & Ordinal & 1,788 & 1,000 \\
RevUtil--Grounding Specificity & Ordinal & 2,652 & 1,000 \\
RevUtil--Helpfulness & Ordinal & 2,279 & 1,000 \\
RevUtil--Verifiability & Ordinal & 1,825 & 788 \\
\hline
\textbf{Total} & -- & \textbf{13,680} & \textbf{5,641} \\
\hline
\end{tabular}
\caption{Task formats and final training and test set sizes. Binary tasks use
scores of 0--1, whereas ordinal tasks use scores of 1--5.}
\label{tab:dataset-process-statistics}
\end{table}
